% Part: proof-theory
% Chapter: natural-deduction
% Section: translation-G2i

\documentclass[../../../include/open-logic-section]{subfiles}

\begin{document}

\olfileid{pt}{nat}{tgi}

\olsection{Nerjemahaké saka \Log{G2i} menyang \Log{N2i}}

Antesèden sekuèn ing \Log{G2i} iku multihimpunan~$\Gamma$
saka !!{formula}s, déné antesèden sekuèn ing \Log{N2i} iku
himpunan~$\Gamma'$ saka formula mawa label, kanthi saben label
katon mung sapisan. Himpunan~$\Gamma'$ saka !!{formula}s mawa label
diarani \emph{cocog karo} multihimpunan~$\Gamma$ yèn, kanggo
saben !!{formula}~$!A$, cacah !!{element}s ing~$\Gamma'$ kang wujudé
$x : !A$ ora ngungkuli multiplisitas~$!A$ (yaiku cacah salinan~$!A$)
ing~$\Gamma$.

\begin{prop}\ollabel{prop:G2i-to-N2i}
Yèn $\Log{G2i} + \Cut \Proves \Gamma \Sequent \Delta$,
banjur $\Log{N2i} \Proves \Gamma' \Sequent !C$,
kanthi $!C \ident \lfalse$ yèn $\Delta$ kosong, déné yèn
$\Delta = \{!A\}$, konklusiné formula tunggal mau, lan
$\Gamma'$ cocog karo~$\Gamma$.
\end{prop}

\begin{proof}
  Kita tuduhaké yèn $\pi$ iku !!a{proof} kanggo
  $\Gamma \Sequent \Delta$ ing \Log{G2i} ditambah aturan~\Cut,
  banjur ana~$\Gamma'$ kang cocog karo~$\Gamma$ lan
  !!a{proof}~$\delta$ kanggo~$\Gamma' \Sequent !C$ ing~\Log{N2i}.
  Iki dibuktèkaké kanthi induksi marang $n = \pheight{\pi}$.

  Yèn $n = 0$, banjur $\pi$ iku aksioma. Yèn iku aksioma
  $\lfalse \Sequent \quad$, $\delta$ iku aksioma
  $x:\lfalse \Sequent \lfalse$, yaiku
  $\Gamma' = \{x:\lfalse\}$ lan $!C \ident \lfalse$.
  Yèn $\pi$ aksioma kang wujudé $!A \Sequent !A$,
  banjur $\delta$ iku $x: !A \Sequent !A$.
  
  Yèn $n>0$, $\pi$ pungkasané aturan~$R$. Hipotesis induksi
  njamin anané !!{proof}s sekuèn kang cocog karo premis-premis
  aturan mau ing~\Log{N2i}. Kita bisa nganggep kabèh label
  pambusakan ing bukti-bukti mau béda-béda, lan label $x$
  mung katon ing sadhuwuré inferènsi kang mbusak label iku,
  yèn pancèn dibusak (kanthi ngowahi label yèn perlu).
  Banjur kita tuduhaké carané nggandhèngaké !!{proof}s mau
  kanggo ngasilaké bukti $\Gamma' \Sequent !C$ kang trep.
  Kasus-kasusé dibédakaké miturut~$R$.

  \begin{enumerate}
    \item Yèn $R$ iku \RightR{\Weakening}, $\pi$ pungkasané
    \[
    \AxiomC{}
    \RightLabel{$\pi_1$}
    \Deduce$\Gamma \fCenter$
    \RightLabel{\RightR{\Weakening}}
    \UnaryInf$\Gamma \fCenter !A$
    \DisplayProof
    \]
    Hipotesis induksi kang ditrapaké marang $\pi_1$ ngasilaké
    !!a{proof} kanggo~$\Gamma' \Sequent \lfalse$.
    Kita ngetrapaké $\FalseInt$ kanggo ngasilaké
    !!a{proof} kanggo $\Gamma' \Sequent !A$.
    \item Yèn $R$ iku \LeftR{\Weakening}, $\pi$ pungkasané
    \[
    \AxiomC{}
    \RightLabel{$\pi_1$}
    \Deduce$\Gamma \fCenter \Delta$
    \RightLabel{\LeftR{\Weakening}}
    \UnaryInf$!A, \Gamma \fCenter \Delta$
    \DisplayProof
    \]
    Hipotesis induksi kanggo $\pi_1$ ngasilaké !!a{proof}
    kanggo~$\Gamma' \Sequent !C$; iki dipilih dadi $\delta$.
    Élinga yèn $\Gamma'$ cocog karo $\Gamma$, uga cocog karo
    $!A, \Gamma$, amarga cacah $x : !A$ ing $\Gamma'$ ora ngungkuli
    cacah salinan~$!A$ ing $\Gamma$, mula uga ora ngungkuli
    cacah salinan $!A$ ing $\Gamma \cup \{!A\}$.
    \item Yèn $R$ iku \LeftR{\Contraction}, $\pi$ pungkasané
    \[
    \AxiomC{}
    \RightLabel{$\pi_1$}
    \Deduce$!A, !A, \Gamma \fCenter \Delta$
    \RightLabel{\LeftR{\Contraction}}
    \UnaryInf$!A, \Gamma \fCenter \Delta$
    \DisplayProof
    \]
    Hipotesis induksi kanggo $\pi_1$ ngasilaké !!a{proof}~$\delta_1$
    kanggo~$\Pi, \Gamma' \Sequent !C$, kanthi $\Pi \subseteq \{x : !A,
    y:!A\}$. Yèn $\Pi = \{x : !A, y:!A\}$, ganti kabèh
    !!{formula}s mawa label $y:!A$ ing~$\delta_1$ dadi $x:!A$.
    Amarga $x$~katon ing kontèks sekuèn pungkasan saka~$\delta_1$,
    kita bisa nganggep ora ana inferènsi ing~$\delta_1$ kang
    mbusak !!a{formula} kang wujudé~$x:!B$; tegesé, ora ana
    $x:!B$ ing sisih kiwa sekuèn ing~$\delta_1$ kang aktif.
    Mula asilé isih !!a{proof} ing~\Log{N2i}.
    \item Yèn $R$ iku $\RightR{\lif}$, $\pi$ pungkasané
    \[
    \AxiomC{}
    \RightLabel{$\pi_1$}
    \Deduce$!A, \Gamma \fCenter !B$
    \RightLabel{\RightR{\lif}}
    \UnaryInf$\Gamma \fCenter !A \lif !B$
    \DisplayProof
    \]
    Miturut hipotesis induksi, ana !!a{proof}~$\delta_1$ kanggo $x:
    !A, \Gamma' \Sequent !B$ utawa $\Gamma' \Sequent !B$,
    kanthi $\Gamma'$ cocog karo~$\Gamma$. Bukti~$\delta_1$
    bisa ditambahi kanggo ngasilaké~$\delta$ kanthi panrapan
    \Intro{\lif}.
    \item Yèn $R$ iku \LeftR{\lif}, $\pi$ pungkasané
    \[
    \AxiomC{}
    \RightLabel{$\pi_1$}
    \Deduce$\Gamma_1 \fCenter !A$
    \AxiomC{}
    \RightLabel{$\pi_2$}
    \Deduce$!B, \Gamma_2 \fCenter \Delta$
    \RightLabel{\LeftR{\lif}}
    \BinaryInf$!A \lif !B, \Gamma_1, \Gamma_2 \fCenter \Delta$
    \DisplayProof
    \]
    Hipotesis induksi ngasilaké !!{proof}s $\delta_1$ kanggo
    $\Gamma_1' \Sequent !A$ lan $\delta_2$ kanggo
    $x: !B, \Gamma_2' \Sequent !C$ utawa kanggo
    $\Gamma_2' \Sequent !C$. Yèn $\delta_2$ wujudé kang kapindho,
    $\delta_2$ bisa dadi~$\delta$. Yèn ora, kita bisa ngowahi
    label formula lan label pambusakan ing~$\delta_1$ supaya
    ora ana label kang katon ing $\delta_1$ lan~$\delta_2$
    bebarengan. Banjur $\Gamma_1' \cup \Gamma_2'$ cocog karo
    $\Gamma_1 \cup \Gamma_2$. Amarga $x:!B$ katon ing sekuèn
    pungkasan saka~$\delta_2$, $x:!B$ ora bisa aktif ing inferènsi
    mawa label pambusakan~$x$ ing~$\delta_2$. Ganti saben aksioma
    $x:!B \Sequent !B$ ing~$\delta_2$ nganggo !!{proof}~$\delta_3$
    \[
    \Axiom$x: !A \lif !B \fCenter !A \lif !B$
    \AxiomC{}
    \RightLabel{$\delta_1$}
    \Deduce$\Gamma_1' \fCenter !A$
    \RightLabel{\Elim{\lif}}
    \BinaryInf$x: !A \lif !B \fCenter !B$
    \DisplayProof
    \]
    sarta saben kemunculan $x:!B$ ing~$\delta_2$ nganggo
    $x:!A \lif !B$. Asilé $\Subst{\delta_2}{\delta_3}{x:!B}$
    iku !!a{proof} kanggo
    $x:!A \lif !B, \Gamma_1', \Gamma_2' \Sequent !C$.
    \item Yèn $R$ iku \RightR{\land}, $\pi$ pungkasané
    \[
    \AxiomC{}
    \RightLabel{$\pi_1$}
    \Deduce$\Gamma_1 \fCenter !A$
    \AxiomC{}
    \RightLabel{$\pi_2$}
    \Deduce$\Gamma_2 \fCenter !B$
    \RightLabel{\RightR{\land}}
    \BinaryInf$\Gamma_1, \Gamma_2 \fCenter !A \land !B$
    \DisplayProof
    \]
    Hipotesis induksi ngasilaké bukti $\delta_1$ kanggo
    $\Gamma_1' \Sequent !A$ lan $\delta_2$ kanggo
    $\Gamma_2' \Sequent !B$. Kanthi ngowahi label formula
    lan label pambusakan ing~$\delta_2$, kita bisa njamin ora ana
    label kang katon ing $\delta_1$ lan $\delta_2$ bebarengan.
    Banjur $\Gamma_1' \cup \Gamma_2'$ cocog karo
    $\Gamma_1 \cup \Gamma_2$. Kita olèh !!a{proof}~$\delta$
    kanthi ngetrapaké $\Intro{\land}$ marang sekuèn pungkasan
    saka $\delta_1$ lan~$\delta_2$.
    \item Yèn $R$ iku $\LeftR{\land}$, $\pi$ pungkasané,
    tuladhané,
    \[
    \AxiomC{}
    \RightLabel{$\pi_1$}
    \Deduce$!A, \Gamma \fCenter \Delta$
    \RightLabel{\LeftR{\land}}
    \UnaryInf$!A \land !B, \Gamma \fCenter \Delta$
    \DisplayProof
    \]
    Miturut hipotesis induksi ana !!a{proof}~$\delta_1$ kanggo
    $x: !A, \Gamma' \Sequent !C$ utawa $\Gamma' \Sequent !C$,
    kanthi $\Gamma'$ cocog karo~$\Gamma$. Ing kasus kapindho,
    bukti~$\delta_1$ bisa langsung dadi~$\delta$. Ing kasus kapisan,
    amarga $x:!A$ katon ing sekuèn pungkasan, $x:!A$ ora aktif
    ing inferènsi mawa label pambusakan~$x$ ing~$\delta_1$.
    Ganti saben aksioma $x:!A \Sequent !A$ nganggo
    !!{proof}~$\delta_2$
    \[
    \Axiom$x: !A \land !B \fCenter !A \land !B$
    \RightLabel{\Elim{\land}}
    \UnaryInf$x: !A \land !B \fCenter !A$
    \DisplayProof
    \]
    lan ganti saben kemunculan $x:!A$ dadi $x:!A \land !B$,
    yaiku mbentuk $\Subst{\delta_1}{\delta_2}{x:!A}$.
    Iki !!a{proof} kanggo $x:!A \land !B, \Gamma' \Sequent !C$.
    \item Yèn $R$ iku \Cut, $\pi$ pungkasané
    \[
    \AxiomC{}
    \RightLabel{$\pi_1$}
    \Deduce$\Gamma_1 \fCenter !A$
    \AxiomC{}
    \RightLabel{$\pi_2$}
    \Deduce$!A, \Gamma_2 \fCenter \Delta$
    \RightLabel{\Cut}
    \BinaryInf$\Gamma_1, \Gamma_2 \fCenter \Delta$
    \DisplayProof
    \]
    Hipotesis induksi ngasilaké bukti $\delta_1$ kanggo
    $\Gamma_1' \Sequent !A$ lan $\delta_2$ kanggo
    $x:!A, \Gamma_2' \Sequent !C$ utawa kanggo
    $\Gamma_2' \Sequent !C$. Ing kasus kapindho,
    pilih $\delta$ dadi~$\delta_2$. Yèn ora, ganti saben aksioma
    $x:!A \Sequent !A$ ing $\delta_2$ nganggo $\delta_1$
    supaya ngolèhké~$\delta$ minangka !!{proof}~$\Subst{\delta_2}{\delta_1}{x:!A}$
    kanggo $\Gamma_1', \Gamma_2' \Sequent !C$.
  \end{enumerate}
\end{proof}

\end{document}
